% Solutions/hints for ch07 exercises (assembled into Appendix C)

\begin{solution}{exr:ch07-classify}
(a) $\pi_1[1]=\pi_2[0]=B$: a following conflict only; allowed. (b) $\pi_1[1]=\pi_2[0]$ and $\pi_2[1]=\pi_1[0]$: a swapping conflict, that is, the degenerate case $m=2$ of the rotation pattern (which is why cycle conflicts are defined for $m\ge3$), and a pair of following conflicts; forbidden. (c) $\pi_1[1]=\pi_2[1]=B$: a vertex conflict at $t=1$; forbidden. (d) Both agents traverse $A\to B$: an edge conflict, which by \cref{prop:ch07-hierarchy}(a) contains the vertex conflicts at $A$ ($t=0$) and at $B$ ($t=1$); forbidden. (e) Agent $2$ waits at $B$ while agent $1$ enters it: a following conflict that is also the vertex conflict $\langle a_1,a_2,B,1\rangle$ by \cref{prop:ch07-hierarchy}(d); forbidden.
\end{solution}

\begin{solution}{exr:ch07-costs}
$\pi_1$ arrives at its goal $(0,2)$ at $t=3$ and never leaves, so $\cost(\pi_1)=3$: the wait at $(0,1)$ before the arrival is paid, the two waits at the goal are free. $\pi_2$ starts at its goal and $\cost(\pi_2)=0$. $\pi_3$ reaches $(1,1)$ at $t=1$, leaves it at $t=2$ and returns at $t=3$; the last arrival after which the agent never leaves is at $t=3$, so $\cost(\pi_3)=3$ by \cref{eq:ch07-path-cost}. Hence $\sumcost=3+0+3=6$ and $\makespan=3$. $\pi_3$ does not cost $1$ because under stay-at-target the agent is only ``done'' once it stays at its goal forever, and an agent that leaves the goal was not done.
\end{solution}

\begin{solution}{exr:ch07-conflict-cases}
Hints. (a) Two agents crossing an empty $3\times3$ grid through its centre, one horizontally and one vertically, meet at the centre at $t=1$ and nowhere else. (b) Put the agents at the two ends of a corridor of \emph{even} length and let them exchange ends: for four cells $(0,0),\dots,(0,3)$ the positions at $t=1$ are $(0,1)$ and $(0,2)$, at $t=2$ they are $(0,2)$ and $(0,1)$, so the agents swap along the middle edge without ever sharing a cell. In a corridor of odd length the two agents meet in the middle cell instead, which is a vertex conflict and not a swap. (c) A convoy: agent $1$ from $(0,0)$ to $(0,3)$ and agent $2$ from $(0,1)$ to $(0,4)$ on a single row; agent $2$ always leaves the cell that agent $1$ enters. (d) Agent $1$ from $(0,1)$ to $(0,2)$ (cost $1$) and agent $2$ from $(0,0)$ to $(0,3)$: agent $2$ reaches $(0,2)$ at $t=2$ while agent $1$ is parked there. (e) Four agents on a $2\times2$ block, each with its goal at the next cell clockwise, all move at $t=0$; \code{all\_conflicts} returns the empty list because a rotation is allowed, and the plan is valid. All five cases appear in the self-test of \code{ch07\_mapf.py}.
\end{solution}

\begin{solution}{exr:ch07-separation}
At an integer time the two drones sit at distinct lattice vertices, and distinct vertices of a lattice with cell side $\ell$ are at least $\ell$ apart. During a step, each drone moves along a segment of length $\ell$ or waits, and the squared distance between two points moving at constant velocity is a convex quadratic in the fraction $\tau\in[0,1]$ of the step, so the minimum is attained at an endpoint or at a single interior point. Enumerate the pairs of moves that a conflict-free plan allows: two waits, a wait and a move, two parallel moves, two moves on edges without a common vertex, and following moves along the same line or around a corner (moving towards the same vertex is a vertex conflict, and reversing the same edge is a swap). All cases except the last keep the distance at $\ell$ or more, because the segments are parallel or separated by a full cell. For the corner case take drone $1$ moving from $(0,0)$ to $(\ell,0)$ and drone $2$ leaving $(\ell,0)$ for $(\ell,\ell)$ at the same time; their positions are $(\tau\ell,0)$ and $(\ell,\tau\ell)$, the squared distance is $\ell^2\bigl((1-\tau)^2+\tau^2\bigr)$, and its minimum at $\tau=\tfrac12$ is $\ell^2/2$, so the distance is $\ell/\sqrt2$. To keep two drones of diameter $d$ at least $m$ apart at all times one therefore needs $\ell/\sqrt2\ge d+m$, that is $\ell\ge\sqrt2\,(d+m)$; forbidding following conflicts removes the corner case and $\ell\ge d+m$ suffices.
\end{solution}

\begin{solution}{exr:ch07-objectives}
(b) Write $L$ for the common shortest-path length. By \cref{prop:ch07-bounds} every solution has $\cost(\pi_i)\ge \dist(s_i,g_i)=L$ for every $i$, hence $\sumcost\ge kL$ and $\makespan\ge L$. A solution $\Pi^\ast$ with $\makespan(\Pi^\ast)=L$ has $\cost(\pi_i)\le L$ for every $i$ and therefore $\cost(\pi_i)=L$ for every $i$, so $\sumcost(\Pi^\ast)=kL$: both lower bounds are attained, and the optima are $\sumcost^\ast=kL$ and $\makespan^\ast=L$. Now for any solution $\Pi$, $\sumcost(\Pi)=kL$ forces $\cost(\pi_i)=L$ for all $i$ (each summand is at least $L$), which is exactly $\makespan(\Pi)=L$; and conversely $\makespan(\Pi)=L$ forces $\cost(\pi_i)=L$ for all $i$ and hence $\sumcost(\Pi)=kL$. So $\Pi$ is optimal for one objective if and only if it is optimal for the other. The hypothesis matters: without a solution of makespan $L$ the two optima need not be attained together, which is \cref{ex:ch07-disagree}.
\end{solution}

\begin{solution}{exr:ch07-horizon}
(b) Under disappear-at-target an agent occupies its goal at its arrival time and is gone from the next step on. In the naive plan of \cref{ex:ch07-grid}, $a_3$ arrives at $(0,1)$ at $t=1$, so it is still there at $t=1$ and absent from $t=2$: the vertex conflict of $a_1$ and $a_3$ at $(0,1)$ at $t=1$ remains, the swap of $a_1$ and $a_2$ on the edge $(0,1)$--$(0,2)$ between $t=1$ and $t=2$ remains (both agents are still flying), and the vertex conflict of $a_2$ and $a_3$ at $(0,1)$ at $t=2$ disappears. Two of the three conflicts survive. In the implementation, \MapfPosition must return ``absent'' instead of $g_i$ for $t>T_i$, and every conflict test must skip a pair in which one agent is absent; the horizon argument is unchanged, because after $T$ nobody moves and nobody is present twice. \Cref{eq:ch07-path-cost} becomes the \emph{first} arrival time, $\cost(\pi_i)=\min\set{\tau:\pi_i[\tau]=g_i}$, since a path now ends the moment it reaches the goal and can never leave it again.
\end{solution}

\begin{solution}{exr:ch07-corridor}
(a) Number the cells $0,\dots,n-1$ from left to right and let $p_1[t],p_2[t]$ be the positions of the two agents. Because vertex conflicts are forbidden, $p_1[t]\ne p_2[t]$ at every $t$, so the sign of $D[t]=p_2[t]-p_1[t]$ is defined at every $t$; initially $D[0]=n-1>0$. Each agent moves by at most one cell per step, so $\abs{D[t+1]-D[t]}\le 2$. A change of order means $D[t]\ge 1$ and $D[t+1]\le -1$. If $D[t+1]=0$ the two agents share a cell, a vertex conflict. Otherwise $D[t]=1$ and $D[t+1]=-1$, which forces both agents to move towards each other by one cell, that is $p_1[t+1]=p_2[t]$ and $p_2[t+1]=p_1[t]$: a swapping conflict on the edge between the two adjacent cells. Both are forbidden, so $D[t]>0$ for all $t$ and the left-to-right order never changes. The exchange would need $D=-(n-1)<0$, so the instance has no solution. This is the smallest instance in which feasibility fails for a combinatorial rather than a connectivity reason, and it is why the pocket of part~(b) is needed.

(b) Number the corridor $(0,0),\dots,(0,3)$ and put the pocket at $(1,1)$. Let agent~$1$ step aside into the pocket and come back out behind the other agent: $\pi_1 = \bigl((0,0),(0,1),(1,1),(0,1),(0,2),(0,3)\bigr)$ and $\pi_2 = \bigl((0,3),(0,2),(0,1),(0,0)\bigr)$. The two agents never share a cell and never exchange an edge (between $t=2$ and $t=3$ agent~$1$ leaves the pocket for $(0,1)$ while agent~$2$ leaves $(0,1)$ for $(0,0)$, which is a following conflict, not a swap), so the plan is a solution. Its costs are $\cost(\pi_1)=5$ and $\cost(\pi_2)=3$, hence $\sumcost=8$ and $\makespan=5$. Both are optimal: \code{optimal\_makespan\_bruteforce} and \code{optimal\_soc\_bruteforce} of \code{ch07\_mapf.py} return $5$ and $8$ on this instance in the self-test.

(c) \Cref{thm:ch07-feasible} does not promise that an instance with two empty vertices is solvable; it promises that solvability can be \emph{decided} in polynomial time. The corridor with $n=4$ and $k=2$ satisfies $k \le \abs{V}-2$, the decision procedure applies, and its answer is ``no solution''. Push-and-Rotate (\cref{ch:ch11}) is complete for exactly this class, so it does not loop or time out: it reports that the instance is unsolvable. With the pocket of part~(b) the graph is no longer a path, the order argument of part~(a) breaks, and the same solver returns a plan.
\end{solution}
